\section{Race and Death}

\begin{quotex}
In this short segment on death and race from \emph{Sintesi di dottrina della razza}, \textbf{Julius Evola} returns to traditional doctrines. As expected, he makes clear that personality transcends race, i.e., the person creates race, but the race does not create the person. That is because the person transcends his material and biological conditioning, and has a dual destiny. On the material plane, he continues the race with his first-born son, the ``son of duty''. Yet he also has a supernatural destiny. Hence, at a certain stage of life, a man may abandon, after having fulfilled them, his family duties and dedicate himself to that spiritual destiny.

He condemns biological racism as ``telluric'', since it is restricted to the material conditions of life. Furthermore, such views are associated with the denial of an afterlife. This is the path of darkness, the southern path, unlike the path of light, or northern path, described by Evola. His last line is quite harsh: `` they pass onto others the task for which they were inadequate.''

Those commenters who would dismiss us as just another ``traditionalist'' site saying the same things, are not very careful readers (as already has been demonstrated). On the contrary, few such sites adhere to traditional principles and promote lifestyles in conformity to Tradition, as this chapter demonstrates.
\end{quotex}


On the strength of what should be clear at this point, we would like to pause a moment to make precise the limits of the membership of personality to its race. In this regard, we will immediately point out the view that is unacceptable from the traditional perspective: it is the view of those who have conceived race as a purely biological-human, historical, and, to sum it up, only a terrestrial entity, and who now think that such an entity is the goal of every being who belongs to it, that nothing exists superior to race, since race is the origin of every value and they consider that the idea of the extraterrestrial attainment and purpose of the individual is illusory and detrimental: ``to remain faithful to land and race.''

We have already encountered and criticized this conception more than once. Opposed to it, moreover, we can appeal to the racialist criterion for the evaluation of the ``truth''; according to the various ``races of the spirit'' there are, in a particular way, similar conceptions of race itself---and there is no doubt that the conception we just pointed out can be ``true'' only for a telluric race, only the telluric man is able to support unconditionally such limited horizons. Furthermore in this telluric vision of race the assumption, that we have sketched out, of those ``neopagan'' racialists also comes into play, since for them the only conceivable immortality would be that of the survival in the blood, in terrestrial descendants.

It is true that similar positions, today, are presented to us less on the basis of a theoretical value than of a practical and political value---i.e., they seek to consolidate the unity of the race of the people and to concentrate every spiritual energy of the individual into temporal and historical tasks that this entity has to resolve. But it is likewise true that the ancient Aryan civilizations, even with respect to terrestrial, heroic, and political achievements, had their greatness without feeling the need to appeal to these myths, recognizing instead rather different truths; it is quite evident, in fact, that this view about race recalls the \emph{pitryana}, the ``way of the South'', which we previously mentioned, which is opposed the \emph{devayana}, the ``divine way of the North'', that stood alone as defining the highest Aryan ideal.

And the theory of ``double hereditary'' is also related to such an ideal. Personality is not fully accounted for in historical-biological or horizontal heredity: instead, it appears as a principle that, in order to manifest itself in race (race meant here in the restricted sense), stands in itself beyond race but cannot be exhausted in it. To recognize race---as has been made clear on principle---does not mean to diminish the personality: personality owes the living and articulated material for its specific expression to race and to as much as terrestrial heredity includes, through its self-manifestation and action. However, there is a conditioning that is neither passive nor one-sided in that. Every individual also reacts on race and his heredity, on the basis of his own most intimate nature; he develops the substance in which he is manifested, the form beyond, and it is in that way that interracial differentiation and that diverse purity or completeness of types are realized, things which we already mentioned and to which we will now return, with regard to its social effects: here it is a giving as much as a receiving. At the points where the highest equilibrium and adaptation are reached (equilibrium, according to our tripartite view spirit, soul, body , between the various component of the true race), there is as a culmination beyond which the personality has nowhere to go---it has nowhere to go along a horizontal or earthly line. His work remains and belongs along this line and his child is physiologically his lineage. But personality itself, if it reached such a peak, is ``free'' and can approach a perfection that is now properly supernatural.

This is precisely the most ancient Aryan conception, pertaining to those who do not belong properly to the group of spiritual leaders, conceptions that can also be found in the various views and legends of the Western Middle Ages. The scrupulous observance of the earthly laws of race, caste, etc. is prescribed, i.e., \emph{dharma}, to the point of complete compliance. Such laws also require the assurance of a lineage: life, which is received at birth, must be given back before death, with its own mark, to another being, and it is for this reason that the firstborn son is called the ``son of duty''. After the ``active life'', according to Aryan law, one can retire to an ascetic-contemplative life. And the Iranian-Aryan saying is rather expressive: the true task is not only to first procreate oneself on the horizontal line of terrestrial descendants, but also upwards, in the ascending vertical direction. In Western religion all these views were confused---above all, what is pertinent to the active life was sharply disconnected from the contemplative life and almost always the truly traditional solutions were forgotten, according to which the law, which is not of this earth, prolongs, completes, and strengthens the law of the earth. But even more harmful than such confusions are the ``telluric'' racist views alluded to, in case they should be taken seriously and have a future. According to the traditional teaching of the Aryan peoples, the idea instead became fixed that the destiny and the dignity of the personality is essentially supernatural; this goal, however, acts as the highest motor impulse and the deepest animating force at the heart of the expression that race gives to the personality, therefore simultaneously raises race up to a limit, beyond which, after having left a sign of greatness, the same force is liberated and tends to act so that death is precisely the fulfillment---\emph{telos}---a new birth---the third birth of the Indo-Aryan teaching.

Only the mediocre or failures among beings, i.e., those who are unable to completely fulfill the law and earthly duty, who perhaps think that they do not have an afterlife, that they have for their destiny, the self-dissolution into the confused vitality of race, into the collective and terrestrial substance of blood and heredity, surviving only in that way---in a rather relative sense of the word---to the destruction of their physical individuality, and passing onto others the task for which they were inadequate.


\flrightit{Posted on 2014-03-30 by Aeneas}

